\documentclass[11pt]{article}
\usepackage[T1]{fontenc}
\usepackage[utf8]{inputenc}
\usepackage{lmodern}
\usepackage[ngerman]{babel}
\usepackage{amsmath,amssymb,mathtools,array,booktabs,pdflscape,adjustbox,diagbox}
\usepackage{geometry}
\usepackage{microtype}
\geometry{a4paper,margin=1.45cm}
\setlength{\parindent}{1.25em}
\setlength{\parskip}{0pt}
\makeatletter
\renewcommand\footnotesize{\@setfontsize\footnotesize{7.5}{8.5}\selectfont}
\makeatother
\allowdisplaybreaks
\setlength{\emergencystretch}{10em}
\sloppy
\hbadness=10000
\vbadness=10000
\hfuzz=2pt
\vfuzz=10pt
\raggedbottom

% current section packet macros
\providecommand{\frR}{\mathfrak R}
\providecommand{\frS}{\mathfrak S}
\providecommand{\frT}{\mathfrak T}
\providecommand{\frM}{\mathfrak M}
\providecommand{\frN}{\mathfrak N}
\providecommand{\frA}{\mathfrak A}
\providecommand{\frB}{\mathfrak B}
\providecommand{\frX}{\mathfrak X}
\providecommand{\frY}{\mathfrak Y}
\providecommand{\frZ}{\mathfrak Z}
\providecommand{\frG}{\mathfrak G}
\providecommand{\frH}{\mathfrak H}
\providecommand{\frL}{\mathfrak L}
\providecommand{\frU}{\mathfrak U}
\providecommand{\frD}{\mathfrak D}
\providecommand{\frC}{\mathfrak C}
\providecommand{\frP}{\mathfrak P}
\providecommand{\calR}{\mathcal R}
\providecommand{\calS}{\mathcal S}
\providecommand{\calA}{\mathcal A}
\providecommand{\calB}{\mathcal B}
\providecommand{\calT}{\mathcal T}
\providecommand{\End}{\operatorname{End}}
\providecommand{\Aut}{\operatorname{Aut}}
\providecommand{\Mat}{\operatorname{Mat}}
\providecommand{\rad}{\operatorname{rad}}
\providecommand{\rk}{\operatorname{rk}}
\providecommand{\id}{\operatorname{id}}
\providecommand{\op}{\operatorname{op}}

\begin{document}

\clearpage

\section*{40. Nichtkommutative Algebren}

\begin{center}
\emph{Math. Zs. 37 (1933), 514--541}
\end{center}

Die Hauptsätze der kommutativen Algebra sind bekanntlich in der galoisschen Theorie enthalten, der die Theorie der Abspaltungs- und Zerfällungskörper vorausgeht, d.\,h. der Körper, die ausreichen, damit ein vorgegebenes Polynom einen Linearfaktor abspaltet, beziehungsweise ganz in Linearfaktoren zerfällt.

Die entsprechenden Teile der Algebra entwickele ich hier im Nichtkommutativen, insbesondere im Hyperkomplexen. Und zwar arbeite ich dabei prinzipiell mit nichtkommutativen Methoden, der Darstellung in nichtkommutativen Körpern. Zum Schluß zeige ich, wie die oben genannten Sätze der kommutativen Algebra sich ganz parallel laufend mit Darstellung in kommutativen Körpern begründen lassen.

Die zugrunde liegende Darstellungstheorie, der eine kurze Automorphismen-theorie vorausgeht (§1), bedeutet eine Fortentwicklung der auf der Theorie der Darstellungsmoduln beruhenden. Insbesondere betrachte ich neben der direkten Darstellung die reziproke, wobei sich die eine auf die andere zurückführen läßt, die jetzt vermöge des reziproken Darstellungsmoduls erzeugt wird. Der Vorteil ist, daß der reziproke Darstellungsmodul, also ein Doppelmodul, sich auch auffassen läßt als einseitiger Modul nach einem Erweiterungsring (§2). Dadurch wird die Darstellung in nichtkommutativen Körpern zurückgeführt auf die Idealtheorie im Erweiterungsring; die irreduziblen Darstellungsklassen entsprechen den irreduziblen Idealklassen des Erweiterungsrings, in genauer Analogie zu den Tatsachen, die bei der Darstellung hyperkomplexer Systeme in ihrem kommutativen Koeffizientenbereich für das System selbst gelten (§§3, 4).

Von hier aus ergeben sich die Struktursätze für Matrizenringe über nichtkommutativen Körpern vermöge der Bemerkung (§5), daß jeder Unterring eine Darstellung durch diesen Körper, beziehungsweise eine reziproke Darstellung durch den reziprok isomorphen Körper bedeutet. Dieser reziprok isomorphe Körper bildet ein erstes Analogon zu minimalem Abspaltungs- und Zerfällungskörper im Kommutativen, insofern als er alle reziproken Darstellungen ersten Grades vermittelt und bei endlichem Rang nach dem Zentrum auch eine volle Zerfällung in direkte Summanden vom Rang \(1\). Das ergibt die galoissche Theorie für Körper (§6) und drückt zugleich die Tatsache aus (§7), daß reziprok isomorphe Divisionsalgebren, also Körper endlichen Ranges nach dem Zentrum, inverse Klassen in der R. Brauerschen Gruppe der Algebrenklassen erzeugen.

Eine zweite Begründung der galoisschen Theorie, die allgemeiner für einfache Systeme gilt, ist fast unmittelbare Folge eines Satzes über vertauschbare Unterringe, der selbst fast direkt an die vorangeschickte Automorphismenbetrachtung anknüpft. Bis hierher verlaufen die Betrachtungen rein nichtkommutativ. Aber auch die Fragestellung nach den kommutativen Zerfällungskörpern wird nichtkommutativ behandelt vermöge der Darstellung des Zerfällungskörpers durch die Divisionsalgebra nach der in §5 vorausgeschickten Bemerkung (§7). Den Schluß bildet die Übertragung auf das Kommutative (§8) und die Theorie der Zerfällungskörper für beliebige Systeme (§9).

Auf diese Darstellungstheorie im Kommutativen, und auf die irrationalen Faktorensysteme, hat R. Brauer die Theorie der Zerfällungskörper gegründet. Wegen dieser kommutativen Begründung muß er, ebenso wie später Albert, das Zentrum als vollkommenen Körper annehmen, eine nach der nichtkommutativen Begründung unnötige Einschränkung. Mit derselben Einschränkung, ebenfalls mit Darstellungstheorie im Kommutativen, haben R. Brauer und K. Shoda die Theorie weiter entwickelt.

Schließlich sei erwähnt, daß sich für den Fall der Körper eine kurze Darstellung bei van der Waerden findet, im Anschluß an meine Vorlesung vom Sommer 1928. Die dortigen Vereinfachungen gegenüber dieser Vorlesung habe ich zum Teil in einer zweiten Vorlesung, zum Teil erst hier übernommen und weitergeführt. Der Satz über vertauschbare Unterringe (§5) und die darauf fußende zweite Beweismethode der nichtkommutativen galoisschen Theorie (§6, 1 und 2) ist erst nach der zweiten Vorlesung hinzugekommen. Gewisse Schlußweisen der ersten Vorlesung, die Methode der Durchschnittsbildung von Differenzenidealen, haben, obwohl komplizierter, den Vorteil, sich auf Systeme von unendlichem Rang und auf ganzzahlige Systeme übertragen zu lassen. Für kommutative ganzzahlige Systeme kommt man so zu dem Zusammenhang von Idealdifferentiation und Differente.

Ihre Hauptanwendung findet die Theorie der Zerfällungskörper in der Theorie der verschränkten Produkte und ihrer Faktorensysteme, die selbst wieder die Grundlage zahlentheoretischer Anwendung ergeben; hierauf wird hier nicht mehr eingegangen.

\subsection*{§1. Über Automorphismen, Moduln und Doppelmoduln}

Die Darstellungstheorie auf Grund der Darstellungsmoduln beruht auf der Theorie des Automorphismenrings abelscher Gruppen mit oder ohne Operatoren. Die dabei implizit zugrunde liegenden Schlußweisen, also die Beziehungen zwischen Abbildung und Rechengesetzen, sollen zur Vermeidung von Wiederholungen in einigen einfachen Sätzen formuliert werden.

\paragraph{1. Multiplikative Abbildung. Assoziativgesetz.}
Sei zunächst \(\frG\) eine Gruppe ohne Operatoren, \(\frA\) ihr absoluter Automorphismenbereich, d.\,h. das System aller Homomorphismen von \(\frG\) in sich. \(\frA\) ist multiplikativ abgeschlossen, da das Produkt \( \sigma\tau \) durch
\[
        g(\sigma\tau)=(g\sigma)\tau,
        \qquad g\in\frG,\quad \sigma,\tau\in\frA                         \tag{1}
\]
definiert ist und offenbar dem Assoziativgesetz genügt.

\(\frG\) wird bekanntlich zu einer Gruppe mit Operatoren, wenn eine Menge \(\frB\) von Symbolen \(O,H,\ldots\) gegeben ist, derart daß die Verknüpfungen \(gO,gH,\ldots\) eindeutig definierte Elemente aus \(\frG\) liefern und Homomorphismen von \(\frG\) in sich erzeugen:
\[
        (gh)O=gO\cdot hO .
\]
Ist der Operatorenbereich \(\frB\) multiplikativ abgeschlossen, so ist die eindeutige Abbildung von \(\frB\) auf den zugehörigen Teil des Automorphismenbereichs dann und nur dann multiplikativ-homomorph, wenn die Assoziativrelation
\[
        g(OH)=(gO)H, \qquad g\in\frG,\ O,H\in\frB,                       \tag{1a}
\]
erfüllt ist. Entsprechend erhält man reziproken Homomorphismus, wenn Linksoperatoren vorliegen.

\paragraph{2. Operatorhomomorphe Abbildung.}
Ist \(\frG\) Gruppe mit Operatoren, so wird Operatorhomomorphie definiert durch
\[
        (gO)\sigma=(g\sigma)O,                                          \tag{2}
\]
beziehungsweise bei Linksoperatoren durch
\[
        (Og)\sigma=O(g\sigma).                                          \tag{2*}
\]
Sind zwei Operatorenbereiche \(\frB\) und \(\frC\) gegeben, so erzeugen dann und nur dann die \(\frC\)-Elemente \(\frB\)-Automorphismen und zugleich die \(\frB\)-Elemente \(\frC\)-Automorphismen, wenn die Bereiche kommutativ mit \(\frG\) verbunden sind, beziehungsweise wenn das durchlaufende Assoziativgesetz erfüllt ist:
\[
        (gO)H=(gH)O,                                                    \tag{2a}
\]
beziehungsweise
\[
        (OH)g=O(Hg).                                                    \tag{2a*}
\]

\paragraph{3. Moduln und Doppelmoduln nach Ringen.}
Ein Rechtsmodul \(\frM\) nach einem Ring \(\frR\) ist eine additive abelsche Gruppe mit den Elementen aus \(\frR\) als Rechtsoperatoren, wobei neben der Assoziativrelation (1a) auch die Distributivrelationen erfüllt sind:
\[
        (g+h)\rho=g\rho+h\rho,\qquad
        g(\rho+\sigma)=g\rho+g\sigma.                                  \tag{3a}
\]
Entsprechendes gilt für Linksmoduln.

Unter Doppelmoduln sind zwei Arten zu unterscheiden. Rechtsmoduln nach zwei Ringen \(\frR\) und \(\frS\) heißen Doppelmoduln, wenn außer den für \(\frR\) und \(\frS\) einzeln geltenden Verknüpfungen noch
\[
        (m\rho)\sigma=(m\sigma)\rho
\]
gilt. \(\frR\)-Links-, \(\frS\)-Rechts-Moduln heißen Doppelmoduln, wenn das durchlaufende Assoziativgesetz
\[
        \rho(m\sigma)=(\rho m)\sigma
\]
hinzukommt.

Ist der Operatorenbereich \(\frR\) einer additiv geschriebenen abelschen Gruppe \(\frM\) ein Ring, so ist die Abbildung von \(\frR\) in den absoluten Automorphismenring genau dann ringhomomorph, wenn \(\frM\) Rechtsmodul nach \(\frR\) ist; bei Linksmoduln wird sie reziprok ringhomomorph. Ist \(\frM\) Rechtsmodul nach den Ringen \(\frR,\frS\), so ist \(\frM\) genau dann Doppelmodul, wenn \(\frR\) auf einen Unterring des \(\frS\)-Automorphismenrings und zugleich \(\frS\) auf einen Unterring des \(\frR\)-Automorphismenrings ringhomomorph abbildbar ist. Für \(\frR\)-Links-, \(\frS\)-Rechts-Doppelmoduln ist die Abbildung von \(\frR\) reziprok homomorph und die von \(\frS\) direkt homomorph.

Daraus folgt insbesondere:

1) Ist \(\frR\) ein Ring mit Einselement, so ist \(\frR\) direkt isomorph zu seinem Automorphismenring als \(\frR\)-Linksmodul und reziprok isomorph zu seinem Automorphismenring als \(\frR\)-Rechtsmodul.

2) Ist \(\frR\) voller Matrizenring über einem im allgemeinen nichtkommutativen Körper \(A\),
\[
        \frR=\sum c_{ik}A=\sum A c_{ik},
\]
so wird \(A\) reziprok isomorph dem Automorphismenkörper der einfachen Rechtsideale, direkt isomorph dem Automorphismenkörper der einfachen Linksideale.

\paragraph{4. Übergang von Rechts-Doppelmodul zu einseitigem Modul nach Produktring.}
Auf diesem Übergang beruht wesentlich die Bedeutung des Rechts-Doppelmoduls.

\paragraph{Satz.}
Ist \(\frM\) Rechtsmodul nach einem Ring \(\frT\), der zwei elementweise miteinander vertauschbare Unterringe \(\frR\) und \(\frS\) enthält, so läßt sich \(\frM\) auch als Doppelmodul nach \(\frR\) und \(\frS\) auffassen. Liegt umgekehrt ein Doppelmodul nach \(\frR\) und \(\frS\) vor und existiert ein Produktring \(\frT\), in dem \(\frR\) und \(\frS\) elementweise vertauschbar sind, so läßt sich \(\frM\) auch als Rechtsmodul nach \(\frT\) auffassen, wobei die \(\frT\)-Verknüpfung die gegebene \(\frR\)- und \(\frS\)-Verknüpfung fortsetzt.

Der Produktring im absoluten Automorphismenring besteht aus allen Elementen der Form
\[
        \sum_i \rho_i\sigma_i+\rho+\sigma
\]
(besitzen \(\frR,\frS\) Einselemente, so fallen die Zusatzglieder \(\rho,\sigma\) fort). Die Homomorphismen
\[
        \frR\to\overline{\frR},\qquad \frS\to\overline{\frS}
\]
ergänzen sich durch
\[
        \sum_i \rho_i\sigma_i+\rho+\sigma
        \longmapsto
        \sum_i \overline\rho_i\,\overline\sigma_i+\overline\rho+\overline\sigma
\]
zu einem Homomorphismus \(\frT\to\overline{\frT}\). Somit ist
\[
        m\Bigl(\sum_i \rho_i\sigma_i+\rho+\sigma\Bigr)
        =\sum_i (m\rho_i)\sigma_i+m\rho+m\sigma
\]
eindeutig und definiert einen \(\frT\)-Modul, der den gegebenen \(\frR,\frS\)-Modul umfaßt.

\subsection*{§2. Reziproke und direkte Darstellung}

\paragraph{1. Linearformenmoduln.}
Ein \(\frS\)-Rechtsmodul \(\frM\), wobei \(\frS\) Ring mit Einselement ist, heißt Linearformenmodul in \(\frS\), wenn \(\frM\) direkte Summe von \(n\) eingliedrigen \(\frS\)-Moduln ist:
\[
        \frM=m_1\frS+\cdots+m_n\frS,
\]
derart daß \(m_i\frS\) zu \(\frS\) operatorisomorph ist, also aus \(m_i s=0\) stets \(s=0\) folgt.

\paragraph{Satz 1.}
Ist \(\frM\) Linearformenmodul in \(\frS\), so wird der \(\frS\)-Automorphismenring \(\frU\) von \(\frM\) reziprok isomorph, nicht nur homomorph, dem Ring \(\overline{\frU}\) aller \(n\)-reihigen Matrizen in \(\frS\).

Ein beliebiger \(\frS\)-Automorphismus \(\alpha\) von \(\frM\) ist durch
\[
        m_i\longmapsto \overline m_i=m_i\alpha
\]
vollständig beschrieben, und daraus folgt
\[
        \sum_i m_i s_i\longmapsto
        \sum_i \overline m_i s_i
        =\sum_i (m_i\alpha)s_i
        =\sum_i (m_i s_i)\alpha.
\]
Schreibt man
\[
        (m_1\alpha,\ldots,m_n\alpha)=(m_1,\ldots,m_n)A,
\]
so gibt die Zuordnung \(\alpha\mapsto A\) eine eineindeutige Beziehung; weiter gilt
\[
        \alpha+\beta\mapsto A+B,\qquad
        \alpha\beta\mapsto BA.
\]
Daher ist der Isomorphismus reziprok. Daraus folgt auch:

\paragraph{Satz 1'.}
Der Automorphismenring eines Linearformenmoduls \(n\)-ten Grades in \(\frS\) läßt eine isomorphe, treue reziproke Darstellung durch den vollen Matrizenring der \(n\)-reihigen Matrizen in \(\frS\) zu.

\paragraph{2. Darstellung und Darstellungsmodul.}
Ein Linearformenmodul \(\frM\) in \(\frS\) heißt reziproker Darstellungsmodul von \(\frR\) in \(\frS\), wenn \(\frM\) Doppelmodul nach \(\frR\) und \(\frS\) und zugleich \(\frR\)- und \(\frS\)-Rechtsmodul ist. Er heißt direkter Darstellungsmodul, wenn \(\frM\) Doppelmodul und zugleich \(\frR\)-Links-, \(\frS\)-Rechtsmodul ist.

\paragraph{Satz 2.}
Jeder reziproke beziehungsweise direkte Darstellungsmodul erzeugt eine Klasse äquivalenter reziproker beziehungsweise direkter Darstellungen von \(\frR\) in \(\frS\), und alle reziproken beziehungsweise direkten Darstellungen werden so erzeugt.

Ist \(\frM\) reziproker Darstellungsmodul, so wird \(\frR\) direkt homomorph auf einen Unterring des \(\frS\)-Automorphismenrings von \(\frM\) abbildbar; nach Satz 1' entsteht daraus eine reziproke Darstellung. Umgekehrt liefert eine reziproke Darstellung durch Zusammensetzung mit dem reziproken Isomorphismus in den \(\frS\)-Automorphismenring einen direkten Homomorphismus von \(\frR\), also einen Darstellungsmodul. Für direkte Darstellungsmoduln ist die entsprechende Homomorphie reziprok. Direkte Darstellungen von \(\frR\) sind reziproke Darstellungen des zu \(\frR\) reziproken Ringes, und umgekehrt.

\paragraph{3. Übergang vom reziproken Darstellungsmodul zum Modul nach Erweiterungsring.}
Für Darstellungsmoduln lautet der Übergangssatz:

Ist \(\frM\) Rechtsmodul nach einem Ring \(\frT\), der zwei elementweise vertauschbare Unterringe \(\frR\) und \(\frS\) enthält, und wird \(\frM\) Linearformenmodul nach \(\frS\), so läßt sich \(\frM\) auch als reziproker Darstellungsmodul von \(\frR\) nach \(\frS\) auffassen. Umgekehrt läßt sich ein reziproker Darstellungsmodul, sofern ein Produktring \(\frT\) existiert, als \(\frT\)-Modul auffassen.

Ist \(\frM\) reziproker Darstellungsmodul von \(\frR\) in \(\frS\) und zugleich \(\frT\)-Modul mit \(\frT\) als Produktring, so ist \(\frR,\frS\)-Isomorphismus von \(\frM\) zu einem Modul \(\frN\) gleichbedeutend mit \(\frT\)-Isomorphismus. Jeder Klasse isomorpher \(\frT\)-Moduln entspricht also eine reziproke Darstellungsklasse von \(\frR\) in \(\frS\), und umgekehrt.

\paragraph{Satz über vertauschbare Matrizen.}
Ist \(\frR\mapsto\frR^*\) eine reziproke Darstellung \(n\)-ten Grades von \(\frR\) in \(\frS\), und bedeutet \(\frB^*\) den Ring aller mit \(\frR^*\) elementweise vertauschbaren \(n\)-reihigen Matrizen in \(\frS\), so ergibt \(\frB^*\) eine reziprok-isomorphe Darstellung der \(\frR,\frS\)-Automorphismen des erzeugenden Darstellungsmoduls, also der \(\frS\)-Automorphismen, die zugleich \(\frR\)-Automorphismen sind; falls der Produktring \(\frT\) existiert, sind dies die \(\frT\)-Automorphismen.

Die \(\frR,\frS\)-Automorphismen bedeuten eine Zuordnung \(m_i\mapsto m_i'\), durch welche dieselbe Darstellung entsteht. Die \(m_i'\) brauchen keine \(\frS\)-Basis von \(\frM\) zu bilden; dementsprechend brauchen die Matrizen aus \(\frB^*\) nicht umkehrbar zu sein. In diesem übertragenen Sinn bleibt
\[
        (m_1',m_2',\ldots,m_n')a=(m_1',m_2',\ldots,m_n')A
\]
richtig, ohne daß \(A\) eindeutig durch \(a\) bestimmt ist.

Die Diagonalmatrizen \(E\cdot s\) liefern eine direkt isomorphe Darstellung von \(\frS\), die identische von \(\frS\). Die durch \(\frR\) und \(\frS\) definierte Zuordnung von \(\frT\) zu Matrizen in \(\frS\) ist also keine Darstellung des Ringes \(\frT\), sondern die Erweiterung der reziproken Darstellung von \(\frR\) zu einer nach \(\frS\) operatorhomomorphen:
\[
        \sum_i r_i s_i\longmapsto \sum_i R_i s_i.
\]
Insbesondere ist die reziproke Darstellung von \(\frR\) operatorhomomorph in bezug auf den im Zentrum von \(\frR\) gelegenen Durchschnitt \([\frR,\frS]\) von \(\frR\) und \(\frS\).

\subsection*{§3. Moduln in bezug auf einen Körper}

Im folgenden wird es sich nur um Darstellungen in im allgemeinen nichtkommutativen Körpern handeln. Es sollen daher einige einfache Sätze über Linearformenmoduln nach einem Körper zusammengestellt werden.

\paragraph{1. Normalbasis eines Untermoduls nach gegebener Basis des vollen Moduls.}
Sei \(A\) Körper und
\[
        \frN=x_1A+\cdots+x_nA
\]
Linearformenmodul vom Rang \(n\); weiter sei
\[
        \frL=z_1A+\cdots+z_\ell A
\]
Untermodul vom Rang \(\ell\le n\). Die \(z_i\) heißen Normalbasis nach \(x\), wenn sie von der Form
\[
        z_i=x_i-(x_{\ell+1}\alpha_{i,\ell+1}+\cdots+x_n\alpha_{i,n}),
        \qquad i=1,\ldots,\ell,
\]
sind. Jeder Modul \(\frL\) besitzt nach geeigneter Numerierung der \(x_i\) eine Normalbasis. Denn nach geeigneter Numerierung kommt
\[
        \frN=\frL+x_{\ell+1}A+\cdots+x_nA,
\]
also
\[
        x_i\equiv x_{\ell+1}\alpha_{i,\ell+1}+\cdots+x_n\alpha_{i,n}
        \pmod{\frL},\qquad i=1,\ldots,\ell,
\]
und somit liegen die entsprechenden
\[
        z_i=x_i-(x_{\ell+1}\alpha_{i,\ell+1}+\cdots+x_n\alpha_{i,n})
\]
in \(\frL\) und erschöpfen \(\frL\).

\paragraph{2. Erweiterungsmodul.}
Sei \(A\) Körper, \(P\) Unterkörper. Ein \(A\)-Modul \(N\) vom Rang \(n\) heißt Erweiterungsmodul eines \(P\)-Moduls \(M\), geschrieben \(N=M_A\), wenn \(M\) Untermodul gleichen Ranges von \(N\) ist. Ist
\[
        M=y_1P+\cdots+y_nP,
\]
so wird
\[
        M_A=y_1A+\cdots+y_nA.
\]
Umgekehrt existiert zu jedem \(P\)-Modul \(M\), bis auf Modulisomorphie, eindeutig der Erweiterungsmodul.

\paragraph{Folgerung a).}
Sind \(z_1,\ldots,z_s\) bezüglich \(P\) linear-unabhängige Elemente aus \(M\), so sind sie auch linear-unabhängig nach \(A\) in \(M_A\). Jeder \(P\)-Modul
\[
        T=z_1P+\cdots+z_sP
\]
aus \(M\) erzeugt also einen Erweiterungsmodul
\[
        T_A=z_1A+\cdots+z_sA\subset M_A.
\]

\paragraph{Folgerung b).}
Jeder \(P\)-Modul \(T\) aus \(M\) ist Verengungsmodul, d.\,h. Durchschnitt seines Erweiterungsmoduls mit \(M\):
\[
        T=T_A\cap M .
\]

\paragraph{3. Satz über invariante Moduln.}
Sei \(N=M_A\) Erweiterungsmodul eines \(P\)-Moduls \(M\), und sei \(P\) voller Invariantenkörper gegenüber einer Gruppe \(\frG\) von Ring-Automorphismen von \(A\). Die Gruppe \(\frG\) wird als Operatorenbereich von \(N\) definiert durch
\[
        Gm=m,\qquad m\in M,
\]
und
\[
        G\Bigl(\sum_i m_i\alpha_i\Bigr)=\sum_i m_i\,G(\alpha_i).
\]
Hilfssatz: \(M\) besteht genau aus den bei \(\frG\) ruhenden Elementen von \(N\).

\paragraph{Satz.}
Die Erweiterungsmoduln \(L_A\) der Untermoduln \(L\) von \(M\), und nur diese, sind zulässige Untermoduln gegenüber \(\frG\) als Operatorenbereich.

Die eine Richtung ist klar. Sei umgekehrt \(L\) zulässig gegenüber \(\frG\), und sei \(z_1,\ldots,z_\ell\) eine Normalbasis von \(L\) nach einer \(P\)-Basis \(x_1,\ldots,x_n\) von \(M\). Für jedes \(G\in\frG\) ist dann
\[
        G(z_i)=x_i-(x_{\ell+1}G(\alpha_{i,\ell+1})+\cdots+x_nG(\alpha_{i,n}))
\]
wieder Element von \(L\). Koeffizientenvergleich in \(x_1,\ldots,x_\ell\) ergibt \(G(z_i)=z_i\). Also liegen die \(z_i\) nach dem Hilfssatz in \(M\), und \(L\) ist Erweiterung des \(P\)-Moduls
\[
        T=z_1P+\cdots+z_\ell P
\]
aus \(M\), wobei \(T=L\cap M\).

\subsection*{§4. Hyperkomplexe Systeme und ihre Darstellungsklassen}

\paragraph{1. Erweiterungssatz.}
Wir verbinden die Darstellungssätze von §2 mit den Modulsätzen aus §3. An Stelle der allgemeinen Ringe \(\frR\) betrachten wir hyperkomplexe Systeme mit Einselement \(S,T,\ldots\) über einem kommutativen Körper \(P\):
\[
        S=x_1P+\cdots+x_nP.
\]
An Stelle der beliebigen Darstellungsringe \(\frS\) betrachten wir nur Körper \(A,B,\ldots\), wenn nichts anderes bemerkt wird solche mit Zentrum \(P\). Dann existiert stets der Produktring, in dem \(S\) und \(A\) elementweise vertauschbar sind, mit Durchschnitt \(P\); es handelt sich um das direkte Produkt \(S\times A\), das zugleich Erweiterungsmodul \(S_A\) von \(S\) ist:
\[
        S_A=x_1A+\cdots+x_nA.
\]
Es wird daher auch als Erweiterungsring bezeichnet.

Bei dieser Spezialisierung geht der Satz über invariante Moduln in den Erweiterungssatz über:

\paragraph{Erweiterungssatz.}
Jedes zweiseitige Ideal aus \(S_A\) ist Erweiterungsideal eines Ideals aus \(S\), und zwar Erweiterung seines Durchschnitts mit \(S\).

Ist nämlich die zugrunde gelegte Gruppe \(\frG\) die der inneren Automorphismen von \(A\), so ist \(P\) als Zentrum von \(A\) voller Invariantenkörper, während die elementweise Vertauschbarkeit von \(A\) mit \(S\) aussagt, daß \(\frG\) auf \(S_A\) so erweitert ist, daß \(S\) voller Invariantenbereich wird. Jedes zweiseitige Ideal \(\mathfrak a\) aus \(S_A\) ist zulässige Untergruppe wegen \(\alpha^{-1}\mathfrak a\alpha\subseteq\mathfrak a\), also Erweiterung seines Durchschnittsmoduls \(\mathfrak a\cap S\), der ein Ideal in \(S\) ist.

Zusatz: Ist \(S\) kommutativ, so wird \(S\) Zentrum von \(S_A\). Allgemeiner wird das Zentrum von \(S\) auch Zentrum von \(S_A\). Unter \(A_r,B_n,\ldots\) sollen stets Matrizenringe des angegebenen Grades über \(A,B,\ldots\) verstanden sein:
\[
        A_r=\sum A c_{ik}=\sum c_{ik}A.
\]

Im folgenden wird wesentlich der Erweiterungssatz für einfache Systeme benutzt:

Ist \(S\) einfaches System, so wird auch \(S_A\) einfach:
\[
        S_A=\sum B c_{ik}=\sum c_{ik}B=B_t
\]
mit zugeordnetem Körper \(B\), dessen Zentrum \(P\) umfaßt. Dabei kann \(A\), und damit auch \(B\), von unendlichem Rang über \(P\) sein; nur \(S\) ist als hyperkomplex vorausgesetzt. Besitzen \(S\) und \(A\) das gemeinsame Zentrum \(P\), so wird \(P\) auch Zentrum von \(S_A\). Allgemeiner ist für einfaches hyperkomplexes \(S\) auch das direkte Produkt \(S\times A_r\) einfach.

\paragraph{2. Darstellungsklassen.}
Unter einer reziproken oder direkten Darstellung eines hyperkomplexen Systems verstehen wir eine von der Nulldarstellung verschiedene Darstellung, die operatorhomomorph in bezug auf den Koeffizientenbereich \(P\) ist.

\paragraph{Satz über die Darstellungsklassen.}
Ist \(S\) hyperkomplex mit Einselement und \(P\) als Koeffizientenbereich, \(A\) ein Körper mit Zentrum \(P\), so besitzt \(S\) so viele verschiedene irreduzible reziproke Darstellungsklassen in \(A\), wie es Klassen einfacher Rechtsideale im Restklassenring von \(S_A\) nach seinem Radikal gibt. Ist \(S\) einfaches System, so besitzt es genau eine irreduzible reziproke und eine irreduzible direkte Darstellungsklasse in \(A\).

Nach der Folgerung aus dem Übergangssatz entsprechen die einfachen reziproken Darstellungsklassen und die einfachen Modulklassen nach \(S_A\) einander eineindeutig. Da \(S_A\) von endlichem Rang nach dem Körper \(A\) ist, erfüllt es die Maximal- und Minimalbedingung für einseitige Ideale. Ist \(S\) einfach, so auch \(S_A\); folglich gibt es nur eine einfache einseitige Idealklasse und damit nur eine irreduzible reziproke Darstellungsklasse. Dasselbe gilt reziprok für die direkten Darstellungen.

\paragraph{3. Rangrelation.}
Ist \(S\) einfach und
\[
        S_A=\sum Bc_{ik}
\]
von endlichem Rang sowohl nach \(A\) wie nach \(B\), so ergibt sich
\[
        n=rt,\qquad n=(S:P),\qquad t^2=(S_A:B),                         \tag{1}
\]
wobei \(r\) der Grad der irreduziblen reziproken Darstellung ist.

\paragraph{4. Verschärfung der Rangrelation, wenn \(A\) von endlichem Rang nach \(P\) ist.}
In diesem Fall gelten die drei Relationen
\[
        (S:P)=n=rt,                                                     \tag{1}
\]
\[
        (B:P)\,t=(A:P)\,r,                                               \tag{2}
\]
\[
        (S:P)(B:P)=(A:P)\,r^2.                                           \tag{3}
\]
Dabei bedeutet (2) den Rang nach \(P\) eines einfachen Rechtsideals aus \(S_A\), ausgedrückt durch den Rang nach \(B\) beziehungsweise \(A\); (3) folgt aus (2) durch Multiplikation mit \(r\) vermöge (1).

\subsection*{§5. Anwendung der Darstellungssätze auf Matrizenringe}

Die in §4 betrachtete reziproke Darstellung von \(S\) in \(A\) läßt sich auch auffassen als direkter Homomorphismus von \(S\) in einen Unterring von \(A_r\), wobei \(\overline A\) den zu \(A\) reziprok-isomorphen Körper bezeichnet. Das Prinzip der folgenden Paragraphen ist umgekehrt, die Untersuchung der Matrizenringe \(A_r\) auf die Darstellungstheorie im reziprok-isomorphen Körper \(A\) zurückzuführen.

\paragraph{1. Reduzible und irreduzible Einbettung in \(A_r\).}
Es seien \(A\) und \(\overline A\) reziprok-isomorphe Körper von endlichem oder unendlichem Rang über ihrem Zentrum \(P\). Das einfache System \(S\) über \(P\) heißt irreduzibel beziehungsweise reduzibel in \(A_r\) einbettbar, wenn \(S\) eine irreduzible beziehungsweise reduzible reziproke Darstellung \(r\)-ten Grades in \(A\) besitzt. Ist \(S\) irreduzibel in \(A_r\) einbettbar, so ist es reduzibel in allen Matrizenringen \(A_{rs}\), \(s>1\), einbettbar. Der Isomorphismus ist stets Fortsetzung des identischen von \(P\); nach §4, 3 ist dabei \(r\) ein Teiler des Ranges \(n\) von \(S\) nach \(P\).

Auf diese Weise werden alle \(P\) umfassenden einfachen Unterringe endlichen Ranges nach \(P\) der verschiedenen \(A_r\) erhalten. Denn jeder solche Unterring ist einem Unterring von \(\overline A_r\) reziprok-isomorph und besitzt daher eine reduzible oder irreduzible reziproke Darstellung in \(A\).

\paragraph{2. Satz über innere Automorphismen.}
Sind \(S^{(1)}\) und \(S^{(2)}\) zwei \(P\) umfassende einfache Unterringe von \(A_f\), von endlichem Rang nach \(P\), und sind \(S^{(1)}\) und \(S^{(2)}\) isomorph nach \(P\), so wird dieser Isomorphismus durch einen inneren Automorphismus von \(A_f\) erzeugt.

Denn \(S^{(1)}\) und \(S^{(2)}\) bedeuten, im allgemeinen reduzible, Einbettungen eines einfachen Systems \(S\) in \(A_f\), also direkte Darstellungen \(f\)-ten Grades in \(A\). Sie gehören derselben reduziblen direkten Darstellungsklasse in \(A\) an; die überführende Matrix erzeugt den Automorphismus. Ist \(A\) selbst von endlichem Rang nach \(P\), also hyperkomplex, so geht dies über in den bekannten Satz: Zwei einfache, das Zentrum \(P\) umfassende Untersysteme von \(A_f\), die isomorph nach \(P\) sind, gehen durch einen inneren Automorphismus von \(A_f\) ineinander über. Insbesondere ist jeder Automorphismus von \(A_f\) ein innerer.

\paragraph{3. Satz über elementweise vertauschbare Unterringe.}
Ist \(S\) ein einfaches, \(P\) umfassendes System aus \(A_f\), bedeutet \(R\) die Gesamtheit der mit \(S\) elementweise vertauschbaren Elemente aus \(A_f\), so wird auch \(R\) Matrizenring:
\[
        R=B_s.
\]
Die zugeordneten Körper \(B\) von \(S_A\) und \(B\) von \(R\) werden reziprok isomorph. Der Durchschnitt von \(R\) und \(S\) wird Zentrum von \(S\). \(R\) wird genau dann Körper, wenn \(S\) irreduzibel in \(A_f\) eingebettet ist.

Denn nach §2, 3 wird \(R\) dem Automorphismenring des reziproken Darstellungsmoduls \(\frM\) von \(S\) in \(A\) direkt isomorph; das ist aber der Automorphismenring von \(\frM\), aufgefaßt als \(S_A\)-Modul. Zerfällt \(\frM\) in \(s\) einfache \(S_A\)-Moduln, und schreibt man wie in §4, 1
\[
        S_A=\sum Bc_{ik},
\]
so wird \(R\) isomorph zu \(\sum Bc_{ik}\), wobei die beiden Körper reziprok isomorph sind. \(R=B\) ist genau dann Körper, wenn \(s=1\), also wenn \(S\) irreduzibel eingebettet ist. Der Durchschnitt von \(R\) und \(S\) ist per Definition das Zentrum von \(S\).

\paragraph{4. Verschärfter Vertauschungssatz bei hyperkomplexem Matrizenring.}
In diesem Fall ergeben die Rangrelationen aus §4, 4 die Verschärfung:

Die einfachen, \(P\) umfassenden Untersysteme aus \(A_f\), wobei \(A\) von endlichem Rang nach seinem Zentrum \(P\) ist, zerfallen in Paare \(S,\overline S\), derart daß \(\overline S\) aus der Gesamtheit der mit \(S\) elementweise vertauschbaren Elemente besteht und umgekehrt. Die zugeordneten Körper von \(S_A\) und \(\overline S\) sind reziprok isomorph, ebenso die von \(S\) und \(\overline S_A\). Der Durchschnitt von \(S\) und \(\overline S\) ist das gemeinsame Zentrum. Der eine Unterring ist genau dann Körper, wenn der andere irreduzibel eingebettet ist. Das Produkt der Rangzahlen nach \(P\) von \(S\) und \(\overline S\) ist gleich dem Rang von \(A_f\).

Ist \(f=rs=r's'\), wobei \(S\) irreduzibel in \(A_r\) und \(\overline S\) irreduzibel in \(A_{r'}\) einbettbar ist, so werden die zugeordneten Körper von \(S\) und \(\overline S\) isomorph einem Paar vertauschbarer Unterringe aus \(A_g\) mit \(g=ss'\).

Die Produktaussage folgt unmittelbar aus der Rangrelation. Ist \(S\) irreduzibel eingebettet, so ist \(\overline S\) Körper und reziprok-isomorph zu \(B\), wo \(S_A=B_t\); die Rangrelation (3) aus §4, 4 gibt die Behauptung. Ist \(S\) reduzibel eingebettet, also \(\overline S=B_s\), so ist \(f=rs\), und Multiplikation von (3) mit \(s^2\) ergibt die Aussage. Daraus folgt auch, daß \(\overline S\) genau die Gesamtheit der mit \(S\) elementweise vertauschbaren Elemente ist, und alles Weitere folgt aus dem Satz über vertauschbare Unterringe.

\end{document}
